\documentclass[10pt,a4paper]{amsart}
\usepackage[margin=19mm]{geometry}
\usepackage{amsmath,amssymb,mathtools}
\usepackage[scr=rsfs,cal=euler]{mathalfa}
\usepackage{xfrac}
\usepackage[unicode,hidelinks,bookmarks=false]{hyperref}
\newcommand{\ie}{i.e.\ }
\newcommand{\cf}{cf.\ }
\newcommand{\resp}{respectively\ }
\newcommand{\Subsection}{\subsection}
\providecommand{\nobreakdash}{\nobreak\hbox{-}}
\setlength{\emergencystretch}{2em}
\multlinegap0pt
\usepackage{bm}
\usepackage{bbm}
\usepackage{dsfont}
\usepackage{xspace}
\usepackage{cancel}
\usepackage{mathtools}
\usepackage{cleveref}
\usepackage{amsthm}
\makeatletter
\@ifundefined{theorem}{\newtheorem{theorem}{Theorem}}{}
\@ifundefined{claim}{\newtheorem{claim}[theorem]{Claim}}{}
\@ifundefined{lemma}{\newtheorem{lemma}[theorem]{Lemma}}{}
\makeatother
\makeatletter
\newcommand*{\abs}[1]{\lvert#1\rvert}
\expandafter\ifx\csname poc\endcsname\relax
\newenvironment{poc}{\begin{proof}[Proof of claim]\renewcommand*{\qedsymbol}{$\blacksquare$}}{\end{proof}}
\fi
\makeatother
\title[An edge-spectral Erd\H{o}s-Stone-Simonovits theorem and its ]{An edge-spectral Erd\H{o}s-Stone-Simonovits theorem and its stability}
\author{Yongtao Li, Hong Liu, Shengtong Zhang}
\date{Source: arXiv:2508.15271v1 (CC BY 4.0)}
\begin{document}
\maketitle

\noindent\emph{Source and licence.} An edge-spectral Erd\H{o}s-Stone-Simonovits theorem and its stability. Version arXiv:2508.15271v1 (CC BY 4.0), distributed under the Creative Commons Attribution 4.0 International licence, as recorded in the arXiv licence metadata for this exact version and shown on the abstract page https://arxiv.org/abs/2508.15271v1. Full rights notice: licenses/arxiv-2508.15271-CC-BY-4.0.txt.\par\smallskip

\noindent\textit{Editorial scope.} Counted unit: the Lemma whose source label is \texttt{lem:close-to-Kab} in the article (source0.tex lines 891--899, source label \texttt{lem:close-to-Kab}), delivered together with the article's own complete original proof of it (source0.tex lines 901--1007, one \texttt{proof} environment of 107 lines). No mathematics is written, repaired or re-derived: statement and proof are the article's own text line for line. Required context retained uncounted: lem-square-equal (lines 536--539). Every definition, display and label of those lines is retained unchanged. Omitted from this package and not redistributed: 1--535, 540--890, 900--900, 1008--1753. Nothing inside a retained excerpt is omitted or altered: every retained body is delivered exactly as the article's lines are written, so no omission receipt of any kind accompanies this package. Citations: the retained excerpts carry no citation command at all, so no bibliography entry of the article is transcribed and no reference is redistributed. Reference adaptation: none was needed and none was made; every reference inside the delivered unit resolves inside it, so no unresolved reference or citation is produced. Printed slips kept literal: no printed word, symbol, hypothesis or number of the counted statement or of its proof was altered. Layout and class: the article's own class is \texttt{article} with packages including fancyhdr, babel, amssymb, amsmath, amsthm, amsfonts, enumerate, bm, nicefrac, setspace, scalefnt, xcolor, graphicx, tikz; the wrapper is the lane's pinned \texttt{amsart} editorial wrapper at 10pt on A4, which restates the theorem-like environments the retained text needs and adds the article's own macro definitions that the retained text needs (\texttt{\textbackslash abs}), with extra wrapper packages (bm, bbm, dsfont, xspace, cancel, mathtools, cleveref). The delivered numbering is the unit's own. Assets: no retained span contains a figure environment, a graphics-inclusion command, a PDF-page inclusion, a table or a TikZ picture; no image file is copied into this package, the package declares no asset dependency and no image file is redistributed with it. Licence: version arXiv:2508.15271v1 of this article is distributed under the Creative Commons Attribution 4.0 International licence, as recorded in the arXiv licence metadata for this exact version and shown on the abstract page https://arxiv.org/abs/2508.15271v1. A full rights notice is delivered at licenses/arxiv-2508.15271-CC-BY-4.0.txt. Characters: every retained excerpt is pure ASCII, so the delivered engine, XeLaTeX with its default Latin Modern text font, reports no missing character.
\begin{lemma} \label{lem-square-equal}
  Let $G$ be a bipartite graph with the partition $V(G)=A\sqcup B$, and $\binom{\bm{x}}{ \bm{y}}$ be an eigenvector of $G$ corresponding to $\lambda (G)$ with $\bm{x}$ indexed by $A$ and $\bm{y}$ indexed by $B$. Then we have $\lVert \bm{x}\rVert = \lVert \bm{y}\rVert$, i.e., 
  \[ \sum_{i\in A} x_i^2 = \sum_{j\in B} y_j^2. \]  
\end{lemma}

\begin{lemma}
\label{lem:close-to-Kab}
For every $\varepsilon \in (0,0.01)$, 
let $\delta = \varepsilon^4/100$. 
If $G$ is a bipartite graph with 
$m$ edges and $\lambda (G) \ge 
(1- \delta ) \sqrt{m}$, then there exist disjoint vertex subsets $U, V\subseteq V(G)$ such that 
\[ d(G, K_{U, V}) \leq \varepsilon m. \]  
\end{lemma}

\begin{proof}
 Let $V(G)=A\sqcup B$ be a bipartition of $G$. Without loss of generality, we may assume that $A = \{1,2,\ldots ,a\}$ and $B = \{1,2,\ldots ,b\}$, and the unit Perron--Frobenius eigenvector of $G$ takes values $x_1 \geq \dots \geq x_a$ on $A$ and $y_1 \geq \dots \geq y_b$ on $B$.   
 It follows from Lemma \ref{lem-square-equal} that
    $$\sum_{i \in A} x_i^2 = \sum_{j \in B} y_j^2 = \frac{1}{2}.$$
    Therefore, we have
    $$\sum_{i \in A} \sum_{j \in B} x_i^2y_j^2 = \frac{1}{4}.$$
    On the other hand, the assumption gives 
    $$\sum_{ij \in E(G)} 2x_i y_j = \lambda(G) \geq (1 - \delta) \sqrt{m}.$$
    Combining the two inequalities, we obtain
    \begin{equation}
    \label{eq:Motzkin-Straus-Exploit-I}
    \sum_{ij \notin E(G)} x_i^2 y_j^2 + \sum_{ij \in E(G)} \left(x_i y_j - \frac{1}{2\sqrt{m}}\right)^2 \leq \delta/2,
    \end{equation}
    where $\sum_{ij \notin E(G)}$ goes over each pair $(i, j) \in (A \times B )\backslash E(G)$ exactly once. 

    \medskip 
    Let $r, s$ be the smallest integers such that
   $\sum_{i = 1}^r x_i^2 \geq \varepsilon^2 
    ~\text{and}~  
    \sum_{j = 1}^s y_j^2 \geq \varepsilon^2$, respectively.   

\begin{claim}\label{cl:1}
We have $\frac{1}{3\sqrt{m}} < x_r y_s < \frac{1}{\sqrt{m}}$.
\end{claim}


    \begin{poc}
    It is easy to verify that $(x-\frac{1}{2\sqrt{m}})^2 \ge \frac{1}{4}x^2$ holds for every $x\in (-\infty, \frac{1}{3\sqrt{m}})\cup  (\frac{1}{\sqrt{m}}, +\infty)$. 
    If $x_r y_s \ge \frac{1}{\sqrt{m}}$, then 
    for any $i \le r$ and $j \le s$, we have $x_i y_j \geq \frac{1}{\sqrt{m}}$. This implies
    $$\min\left\{ x_i^2 y_j^2, \left(x_i y_j- \frac{1}{2\sqrt{m}} \right)^2 \right\} \geq \frac{1}{4} x_i^2 y_j^2.$$
    Summing over $i\le r$ and $j\le s$, we have
    $$\sum_{i=1}^r \sum_{j=1}^s \min\left\{ x_i^2 y_j^2, \left(x_i y_j - \frac{1}{2\sqrt{m}}\right)^2 \right\} \geq \frac{1}{4} \left(\sum_{i =1}^r x_i^2 \right) \left( \sum_{j=1}^s y_j^2 \right) \geq \frac{1}{4} \varepsilon^4$$
    which contradicts with  (\ref{eq:Motzkin-Straus-Exploit-I}). If $x_r y_s \leq \frac{1}{4\sqrt{m}}$, then $x_iy_j \le \frac{1}{4\sqrt{m}}$ 
    for $i \ge r$ and $j \ge s$. 
    We can finish by an analogous argument that replaces $i \le r$ and $j \le s$ with $i \geq r$ and $j \geq s$.
    \end{poc}

     Let $U$ and $V$ be sets of vertices of $G$ defined as  
    \[ 
    U := \left\{u\in A: x_u \geq \frac{x_r}{ 4} \right\} 
    \quad~\text{and}~\quad 
    V := \left\{v \in B: y_v \geq \frac{y_s}{4} \right\}. \]
    In the rest of the proof, we use (\ref{eq:Motzkin-Straus-Exploit-I}) to show that $G$ has small edit distance with $K_{U, V}$. 
    
    First of all, we show that there are few  vertex pairs in $U \times V$ that are non-edges in $G$. 

\begin{claim}\label{cl:2}
We have $|E(K_{U, V}) \backslash E(G)| \leq \frac{1}{3} \varepsilon m $.
 \end{claim}  
  
    \begin{poc}
    We need to show that $G$ misses at most $\frac{1}{3} \varepsilon m$ edges between $U$ 
    and $V$.  By~\Cref{cl:1}, for any $(i, j) \in (U\times V) \backslash E(G)$, we have 
    $$x_i y_j \geq \frac{1}{16} x_r y_s  \geq  \frac{1}{48 \sqrt{m}}.$$ 
    Using (\ref{eq:Motzkin-Straus-Exploit-I}), we obtain
        $$|E(K_{U, V}) \backslash E(G)| \leq 
        (48\sqrt{m})^2 \cdot \sum_{ij \notin E(G)} x_i^2 y_j^2
        \leq 
        (48 \sqrt{m})^{2} \cdot \frac{\delta}{2} 
        \leq 1152 \delta m \leq \frac{1}{3} \varepsilon m ,$$
        as desired.
    \end{poc}

    Moreover, we prove that there are few edges in $G[\bar{U},B]$ and 
    $G[A,\bar{V}]$. 

\begin{claim}\label{cl:3}
    We have $e(E(G) \cap (\bar{U} \times B)) \leq \frac{1}{3} \varepsilon m $, and symmetrically $e(E(G) \cap (A \times \bar{V})) \leq \frac{1}{3} \varepsilon m $.
\end{claim} 
    
    \begin{poc}
        Let $E_1$ be the set of edges $ij$ in $E(G) \cap (\Bar{U} \times B)$ with $j \geq s$, and let $E_2$ be the set of edges $ij$ in $E(G) \cap (\Bar{U} \times B)$ with $j < s$. It suffices to show that $\max\{\abs{E_1}, \abs{E_2}\} \leq 
        \frac{1}{6} \varepsilon m $.

        For any $ij \in E_1$, by Claim \ref{cl:1}, we have $x_i y_j \leq \frac{1}{4} x_r y_s  \leq \frac{1}{4\sqrt{m}}$. Then together with (\ref{eq:Motzkin-Straus-Exploit-I}), we have 
        $$\abs{E_1} \cdot \frac{1}{16 m} \leq \sum_{ij \in E_1} \left(x_i y_j - \frac{1}{2\sqrt{m}} \right)^2 \leq \frac{1}{2}\delta , $$
        which implies $\abs{E_1} \leq 8 \delta m < \frac{1}{6} \varepsilon m $, as needed. 

        To bound $|E_2|$, we use the Cauchy--Schwarz inequality to obtain
        \begin{align}
       \notag \sum_{ij \in E(G) \backslash E_2} 2x_i y_j  & \leq 2 \sqrt{\sum_{ij \in E(G) \backslash E_2} x_i^2 y_j^2} \cdot \sqrt{ |E(G) \backslash E_2| } \\ 
        & \leq 2 \sqrt{\sum_{i \in A, j \in B} x_i^2 y_j^2} \cdot \sqrt{|E(G) \backslash E_2|} \notag \\
        &= \sqrt{m - |E_2| }. \label{eq-notin-E2}
        \end{align}
        By the definition of $s$, we have $\sum_{j < s} y_j^2 < \varepsilon^2$. Observe that 
        $$\sum_{ij \in E_2} x_i^2 y_j^2 \leq \left(\sum_{i \in A} x_i^2\right) \left( \sum_{j < s} y_j^2 \right) < \frac{\varepsilon^2}{2}.$$
        Applying the Cauchy--Schwarz inequality, we obtain
        \begin{align} \label{eq-in-E2}
        \sum_{ij \in E_2} 2x_i y_j \leq 2 \sqrt{\sum_{ij \in E_2} x_i^2 y_j^2} \cdot \sqrt{\abs{E_2}} \leq \sqrt{2 \varepsilon^2 \abs{E_2}}.
        \end{align}
        Combining (\ref{eq-notin-E2}) and (\ref{eq-in-E2}), we conclude that 
        $$ 
        (1 - \delta) \sqrt{m} \le 
        \lambda(G) = \sum_{ij \in E(G)} 2x_i y_j 
        \le \sqrt{m - \abs{E_2}} + \sqrt{2 \varepsilon^2 \abs{E_2}}.$$
        By the Cauchy--Schwarz inequality again, we have
        $$\sqrt{m - \abs{E_2}} + \sqrt{2 \varepsilon^2 \abs{E_2}} \leq \sqrt{(1 + 4\varepsilon^2) \left(m - \abs{E_2} + \frac{1}{2}\abs{E_2}\right)}.$$
        Therefore, we obtain
        $$ (1 - \delta)^2 m\le 
        (1 + 4\varepsilon^2) \left(m - \abs{E_2} + \frac{1}{2}\abs{E_2} \right), $$
        which implies that $\abs{E_2} \leq \frac{1}{6} \varepsilon m $ by our choice of $\delta = \varepsilon^4 / 100$ and $\varepsilon \in (0, 0.01)$.
    \end{poc}
    Combining~\Cref{cl:2} and~\Cref{cl:3}, we conclude that 
    $$d(G, K_{U, V}) \leq e(K_{U, V} \backslash G) + e(E(G) \cap (\bar{U} \times B)) + e(E(G) \cap (A \times \bar{V})) \leq \varepsilon m, $$ 
    as desired.
\end{proof}

\end{document}
